\documentclass[a4paper,11pt]{article}
\usepackage{exam-style-842-unit}
\examinehead{842 信息化概论 · 第 3 单元练习册}{练习一\quad 计算机技术基础}
\begin{document}

\unitcover{842 信息化概论 · 单元练习册}{第 3 单元\quad 计算机技术与网络技术}{练习一\quad 计算机技术基础（共 100 分）}{约 70 分钟}{计算机的概念与产生基础；布尔、图灵与冯·诺依曼的贡献；ENIAC 与计算机发展的四个阶段；计算机技术的主要分类（系统技术、器件技术、部件技术、组装技术）；计算机技术的应用；计算机技术的特点；信息技术与计算机技术的关系；计算机技术的主要发展方向与重点发展领域。}

\begin{knowmap}
\textbf{理论准备：}布尔 1854 年建立逻辑代数 $\rightarrow$ 图灵 1936 年提出理想计算机（图灵机）$\rightarrow$ 冯·诺依曼提出存储程序与五大部件 $\rightarrow$ ENIAC 1946 年诞生\\[2pt]
\textbf{冯·诺依曼结构：}运算器 · 控制器 · 存储器 · 输入设备 · 输出设备，核心思想是存储程序\\[2pt]
\textbf{四代计算机：}电子管 $\rightarrow$ 晶体管 $\rightarrow$ 集成电路（三化、小型机）$\rightarrow$ 大规模集成电路（巨型机、微处理器与个人计算机）\\[2pt]
\textbf{四类技术：}系统技术（结构·管理·维护·应用）· 器件技术 · 部件技术（运算控制·存储·输入输出）· 组装技术\\[2pt]
\textbf{技术特点：}高技术五特征（创新·战略·增值·渗透·风险）＋复杂的技术体系·现代科技的技术平台·发展潜力巨大·技术文化已经形成\\[2pt]
\textbf{与信息技术的关系：}计算机技术是信息技术的核心技术之一；二者主要内容、发展历程与应用范围不同\\[2pt]
\textbf{未来发展：}生物·光子·量子·纳米计算机；高性能器件·高性能体系结构·高性能软件·和谐人机交互·纳米技术·网络与多媒体
\end{knowmap}

\begin{keyterms}
\textbf{计算机}：又称电脑，是一种利用电子学原理、根据一系列指令对数据进行处理的工具。\\
\textbf{逻辑代数（布尔代数）}：布尔 1854 年建立，用 0 和 1 表示命题真假并进行代数演算的理论，是计算机逻辑设计的数学基础。\\
\textbf{图灵机}：图灵 1936 年在《理想计算机》论文中提出的通用数字计算机数学模型，为计算机诞生提供了理论准备。\\
\textbf{冯·诺依曼结构}：由运算器、控制器、存储器、输入设备和输出设备五大部分组成、采用存储程序原理的计算机体系结构。\\
\textbf{ENIAC}：1946 年 2 月 14 日诞生于美国宾夕法尼亚大学的第一台电子计算机，为计算导弹弹道而研制。\\
\textbf{计算机系统技术}：把计算机作为一个完整系统所运用的技术，包括系统结构技术、系统管理技术、系统维护技术和系统应用技术。\\
\textbf{计算机器件技术}：以电子器件从真空电子器件到超大规模集成电路的变革推动计算机划代的技术，是计算机系统的物质基础。\\
\textbf{计算机部件技术}：构成计算机系统的各类部件技术，主要有运算与控制技术、信息存储技术和信息输入输出技术。\\
\textbf{计算机组装技术}：解决计算机系统可靠性、维修调试方便性、生产工艺性和信息传递延迟的技术，核心是制冷与提高组装密度。\\
\textbf{操作系统}：管理计算机硬件与软件资源、合理组织计算机工作流程、为用户提供方便接口的最基本的系统软件。\\
\textbf{计算机辅助设计（CAD）}：利用计算机进行产品设计，可提高设计质量、缩短设计周期、提高设计自动化水平。\\
\textbf{实时控制}：又称过程控制，用计算机对连续工作的控制对象及时采集信号、计算处理并发出调节信号进行自动控制。\\
\textbf{模式识别}：用计算机模拟人的智能，对语音、图像、指纹等对象进行辨识的应用。\\
\textbf{计算机技术}：计算机领域中所运用的技术方法和技术手段，具有明显的综合特性，与电子工程、应用物理、机械工程、现代通信技术和数学等紧密结合。
\end{keyterms}

\begin{contrast}{计算机技术 / 信息技术（每行依次为：主要内容·发展历程·应用范围）}
\citem{计算机技术}{硬件、软件及应用三方面的技术，以离散数学、算法论、控制论等为理论基础}{从 ENIAC 至今不到 100 年；若从帕斯卡加法器算起也不到 400 年}{是信息技术应用范围的一部分，以计算机实体为研究对象}
\citem{信息技术}{对信息进行操作的技术，主要支柱是通信技术、计算机技术和控制技术，即“3C”技术}{有信息交换就有信息技术，电报、电话、烽火传书都属于早期信息技术}{外延极大，还包括电视、录音、光盘、网络、数码相机等}
\end{contrast}

\begin{contrast}{系统技术 / 器件技术 / 部件技术 / 组装技术（计算机技术的四个分类；每行依次为：技术定位·主要内容·地位或作用）}
\citem{系统技术}{把计算机作为完整系统所运用的技术}{系统结构技术、系统管理技术、系统维护技术、系统应用技术}{使系统获得良好的解题效率和合理的性能价格比}
\citem{器件技术}{以电子器件的变革推动计算机划代的技术}{从真空电子管到超大规模集成电路}{是计算机系统的物质基础，大大提高性能价格比}
\citem{部件技术}{构成计算机系统的各类部件技术}{运算与控制技术、信息存储技术、信息输入输出技术}{中央处理器处于核心地位，决定运算速度与存储能力}
\citem{组装技术}{解决系统可靠性与组装密度的技术}{制冷与空调、提高组装密度、小型化微型化}{与计算机换代相协调，直接影响可靠性和信息传递延迟}
\end{contrast}

\begin{contrast}{第一代 / 第二代 / 第三代 / 第四代计算机（每行依次为：主要器件·代表成果·主要特点）}
\citem{第一代}{电子管}{ENIAC（1946 年）}{体积大、功耗高、速度慢、可靠性差，主要用于军事和科学计算}
\citem{第二代}{晶体管}{晶体管计算机（1960 年起大量生产）}{引入磁芯存储器与磁盘磁带等外存，速度提高到几十万次每秒，体积功耗下降}
\citem{第三代}{集成电路}{小型计算机}{最重要的特点是实现“三化”：通用化、系列化、标准化}
\citem{第四代}{大规模集成电路}{巨型机、微处理器与个人计算机}{运算与数据处理能力强大，计算机大规模普及}
\end{contrast}

\begin{contrast}{生物 / 光子 / 量子 / 纳米计算机（计算机的主要发展方向；每行依次为：核心思想·主要优势·现状）}
\citem{生物计算机}{以蛋白质分子、DNA 核酶等生物分子作开关和存储}{并行处理能力强，能耗仅为电子计算机的十亿分之一}{受生物技术水平限制，预计 10—20 年后进入实用阶段}
\citem{光子计算机}{以光子代替电子、光互连代替导线互连}{并行处理强、运算速度快、可在室温工作、能耗低}{尚未诞生，处于研究与探索阶段}
\citem{量子计算机}{利用量子叠加与量子并行计算，数据用量子位存储}{存储量大，运算速度可比个人计算机快 10 亿倍}{存在退相干和纠错困难，实验研制已见曙光}
\citem{纳米计算机}{用纳米技术研制芯片，体积仅数百个原子大小}{几乎不耗能、性能强大、生产成本低}{已取得研制进展，是未来芯片制造的重要方向}
\end{contrast}

\qtypename{一、单项选择题}{10 题，每题 3 分，共 30 分}
\q 世界上第一台电子计算机 ENIAC 诞生于（\quad）
\choicesline{A. 1936 年}{B. 1946 年}{C. 1956 年}{D. 1964 年}

\q 提出存储程序原理、其体系结构至今仍为现代计算机普遍采用的是（\quad）
\choicesline{A. 图灵}{B. 布尔}{C. 冯·诺依曼}{D. 巴贝奇}

\q 1936 年发表《理想计算机》论文、提出通用数字计算机数学模型的是（\quad）
\choicesline{A. 图灵}{B. 冯·诺依曼}{C. 布尔}{D. 艾肯}

\q 1854 年建立逻辑代数、为计算机逻辑设计奠定数学基础的是（\quad）
\choicesline{A. 莱布尼茨}{B. 布尔}{C. 帕斯卡}{D. 莫尔斯}

\q 第一代电子计算机采用的主要电子器件是（\quad）
\choicesline{A. 电子管}{B. 晶体管}{C. 集成电路}{D. 大规模集成电路}

\q 第三代计算机最重要的特点是实现了（\quad）
\choices{A. 通用化、系列化、标准化}{B. 微型化、网络化、智能化}
{C. 数字化、多媒体化、网络化}{D. 集成化、高速化、廉价化}

\q 运算与控制技术、信息存储技术和信息输入输出技术属于计算机技术中的（\quad）
\choicesline{A. 系统技术}{B. 器件技术}{C. 部件技术}{D. 组装技术}

\q 计算机最基层的逻辑线路是（\quad）
\choices{A. 与门、或门和反相器}{B. 加法器、寄存器和计数器}
{C. 触发器、译码器和存储器}{D. 组合逻辑和时序逻辑}

\q 利用计算机来进行产品设计、可提高设计质量和缩短设计周期的技术是（\quad）
\choicesline{A. CAM}{B. CAD}{C. CAI}{D. CAT}

\q 下列不属于计算机技术“高技术”特征的是（\quad）
\choicesline{A. 高度的创新性}{B. 高度的战略性}{C. 高度的增值性}{D. 高度的消耗性}

\qtypename{二、多项选择题}{4 题，每题 4 分，共 16 分；每题有两个或两个以上正确项}
\q 操作系统的主要功能包括（\quad）
\choicesline{A. 处理机管理}{B. 存储管理}{C. 文件管理}{D. 设备管理}

\q 计算机技术可粗分为（\quad）
\choicesline{A. 系统技术}{B. 器件技术}{C. 部件技术}{D. 网络技术}

\q 关于信息技术与计算机技术的联系，下列说法正确的有（\quad）
\choices{A. 计算机技术是信息技术的核心技术之一}{B. 计算机是信息接收、处理、存储的主要载体}
{C. 二者都是高技术，都具有高技术的特征}{D. 二者主要内容完全相同}

\q 计算机技术的主要应用领域包括（\quad）
\choicesline{A. 数值计算}{B. 数据处理}{C. 实时控制}{D. 模式识别}

\qtypename{三、判断题}{6 题，每题 2 分，共 12 分；正确的画“√”，错误的画“×”并写出改正后的表述}
\q ENIAC 是世界上第一台电子计算机，1946 年诞生于美国宾夕法尼亚大学。\quad（\quad）\quad 改正：\fixblank

\q 冯·诺依曼结构的计算机采用存储程序原理，由运算器、控制器、存储器、输入设备和输出设备五大部分组成。\quad（\quad）\quad 改正：\fixblank

\q 计算机发展的第四代以晶体管为主要器件。\quad（\quad）\quad 改正：\fixblank

\q 操作系统是最基本的系统软件，它管理计算机系统的硬件与软件资源。\quad（\quad）\quad 改正：\fixblank

\q 计算机技术与信息技术的主要内容、发展历程和应用范围完全相同。\quad（\quad）\quad 改正：\fixblank

\q 计算机技术只包括硬件技术，不包括软件技术。\quad（\quad）\quad 改正：\fixblank

\qtypename{四、填空题}{6 空，每空 1 分，共 6 分}
\q 1946 年诞生于美国宾夕法尼亚大学的第一台电子计算机名称是\fillblanks[2.0cm]。

\q 现代计算机普遍采用的体系结构称为\fillblanks[2.6cm]。

\q 计算机技术可粗分为系统技术、器件技术、部件技术和\fillblanks[2.0cm]。

\q 计算机最基层的逻辑线路为“与门”“或门”和\fillblanks[1.8cm]。

\q 第三代计算机最重要的特点是实现了通用化、系列化和\fillblanks[1.6cm]。

\q 用计算机对连续工作的控制对象实行自动控制的技术称为\fillblanks[2.0cm]。

\newpage
\qtypename{五、简答题}{3 题，每题 8 分，共 24 分}
\q 简述计算机发展所经历的四个阶段（代）及各阶段的主要特征。\anslines{2.4cm}

\q 简述计算机技术的主要分类，并说明各类技术的主要内容。\anslines{2.4cm}

\q 简述信息技术与计算机技术的联系与区别。\anslines{2.4cm}

\qtypename{六、综合题}{1 题，共 12 分}
\q 计算机技术正面临一系列新的重大变革，科学家一方面探索生物、光子、量子、纳米等新型计算机，另一方面大力推进现有计算机技术的重点发展领域。
请回答：
（1）说明计算机技术的主要发展方向，并简述生物计算机、光子计算机、量子计算机和纳米计算机各自的核心思想与主要优势；
（2）说明当前计算机技术的重点发展领域；
（3）结合农业信息化，谈谈上述技术方向可能带来的影响。\anslines{5cm}

\end{document}
